\documentclass[11pt]{article}
\usepackage[margin=1in]{geometry}
\usepackage[T1]{fontenc}
\usepackage{lmodern,amsmath,amssymb,amsthm,mathtools,booktabs,xcolor}
\usepackage{hyperref}
\hypersetup{colorlinks=true,linkcolor=blue!45!black,citecolor=blue!45!black,urlcolor=blue!45!black,
pdfauthor={Joshua Mark},pdftitle={Parallel progress toward exact Hadamard existence}}
\setlength{\emergencystretch}{3em}
\newtheorem{proposition}{Proposition}[section]
\newtheorem{theorem}[proposition]{Theorem}
\newtheorem{corollary}[proposition]{Corollary}
\newcommand{\Z}{\mathbb Z}
\newcommand{\R}{\mathbb R}
\newcommand{\one}{\mathbf 1}
\title{Parallel progress toward exact Hadamard existence}
\author{Joshua Mark\thanks{ORCID: \href{https://orcid.org/0009-0003-5905-2382}{0009-0003-5905-2382}. AI tools assisted with parallel mathematical research, exposition and exact computation.}}
\date{20 November 2026 --- research progress note}
\begin{document}
\maketitle
\begin{abstract}
We make progress on four independent routes to the Hadamard conjecture.
A classical three-square theorem solves a uniform high-energy region of
the cyclic Hermitian norm problem, with a smaller sufficient threshold
at skew-character lengths. Integer moment analysis shows that coordinate
cube moves cannot leave any saturated Hadamard array; richer trades are
necessary. An exact order-twenty-four certificate disproves universal
two-row augmentation. A tensor construction supplies non-Hadamard strict
local minima of the quartic flattening objective at infinitely many
admissible orders. The all-order Hadamard conjecture remains unproved.
These deductions and finite certificates carry no claim of bibliographic
novelty, and none is presented as a Lean formalization.
\end{abstract}

\section{What is established and what is still missing}
The preceding research notebook posed a broader arithmetic question:
for odd $m,t$, can four coefficientwise odd elements of
$\Z[z]/(z^m-1)$ have Hermitian norm sum $4t$? Necessarily $t\ge m$.
At $t=m$, energy forces every coefficient to be a sign, and the
Goethals--Seidel transfer gives order $4m$; doubling then covers all
multiples of four if that boundary can be reached for every odd $m$.
We do not reach that boundary uniformly here.

The main positive results are the uniform sufficient thresholds
\[
 t\ge\frac{m^2+3}{4},\qquad
 t\ge\frac{m^2+6m+21}{16}
 \quad\text{when an appropriate skew character exists}.
\]
The second threshold applies to primes $m\equiv3\pmod4$. Both remain
quadratic rather than linear. They are sharp within the displayed
construction families, not asserted to be optimal for the unrestricted
norm problem. The moment branch also proves an explicit cofinite
run-count bound in each fixed coordinate dimension.

The arithmetic objects also have a direct lattice interpretation.
The same Goethals--Seidel block identity works for integral circulants:
a norm solution gives an odd integral matrix $K$ of order $4m$ with
$KK^{\mathsf T}=4tI$. Each row has $4m$ odd entries, so its squared
length is at least $4m$. Equality is exactly the sign-matrix case.
Thus the new construction gives integral scaled isometries above their
minimum possible energy; it does not attain that minimum uniformly.

The negative results identify concrete missing ingredients. Ordinary
cube-kernel generators are not a positive-conversion mechanism at the
target; two-row exchanges do not always augment a partial matrix; and
even a smooth quartic objective has strict nonglobal minima. These are
obstructions to specific methods, not to the Hadamard conjecture.
All new work remains separate from the existing Lean repository and
the earlier asymptotic manuscript.

\section{Arithmetic: a proved high-energy norm theorem}


Put $R_m=\Z[z]/(z^m-1)$, let $z^*=z^{-1}$, and write
$J=1+z+\cdots+z^{m-1}$. We seek
\begin{equation}\label{eq:target}
 A_i\equiv J\pmod{2R_m},\qquad
 \sum_{i=0}^3A_iA_i^*=4t\cdot1.
\end{equation}
When $m$ is odd, elementary energy and augmentation arguments force
$t\ge m$ and $t$ odd. At $t=m$, every coefficient is a sign, giving
four periodic complementary binary sequences and hence, by the
Goethals--Seidel construction, a Hadamard matrix of order $4m$.
The general boundary problem is not solved here.

\begin{theorem}[Uniform large-energy existence]\label{thm:main}
If $m\ge1$ and $t$ are odd integers and $t\ge(m^2+3)/4$, then
\eqref{eq:target} has a solution of the form
\[
 A_i=J+2x_i\cdot1,\qquad x_i\in\Z.
\]
For odd $m,t$, this particular form exists if and only if
$t\ge(m^2+3)/4$.
\end{theorem}
\begin{proof}
Since $JJ^*=mJ$ and $J^*=J$,
\begin{equation}\label{eq:reduction}
 \sum_{i=0}^3(J+2x_i)(J+2x_i)^*
 =4\Bigl(m+\sum_i x_i\Bigr)J+4\sum_i x_i^2\cdot1.
\end{equation}
For $m>1$, $J$ and $1$ are linearly independent. Thus the target is
equivalent to
\begin{equation}\label{eq:four}
 \sum_i x_i=-m,\qquad \sum_i x_i^2=t.
\end{equation}
For $m=1$, these conditions are still sufficient; necessity of the
stated threshold merely says $t\ge1$ and follows from coefficient energy.

Set $N=4t-m^2$. Under the hypotheses, $N\ge3$ and $N\equiv3\pmod8$.
The Legendre--Gauss three-square theorem gives
$N=u^2+v^2+w^2$ for integers $u,v,w$. All three are odd, because the
only square residues modulo eight are $0,1,4$. By changing the sign of
one of them if necessary, arrange
$u+v+w\equiv m\pmod4$. Define
\begin{align*}
 x_0&=(-m+u+v+w)/4,& x_1&=(-m+u-v-w)/4,\\
 x_2&=(-m-u+v-w)/4,& x_3&=(-m-u-v+w)/4.
\end{align*}
They are integers: the last three numerators differ from the first by
twice a sum of two odd numbers. Their sum is $-m$, and the orthogonal
four-square identity gives
\[
 \sum_i x_i^2=\frac{m^2+u^2+v^2+w^2}{4}=t.
\]
Equation~\eqref{eq:reduction} proves existence.

Conversely, for $m>1$ any solution of this form obeys
\eqref{eq:four}. Cauchy--Schwarz gives $4t\ge m^2$. Since
$4t-m^2\equiv3\pmod8$, it is at least three. This is precisely
$t\ge(m^2+3)/4$.
\end{proof}

The invoked classical three-square theorem has an elementary geometric
proof in Ankeny~\cite{Ankeny}, whose Theorem 1 applies whenever the
integer is not $4^a(8b+7)$. No code or proof supplement from that paper
is used.

\begin{corollary}
The broader parity-constrained norm problem has a solution for every
odd $t\ge m$ when $m=1$ or $m=3$.
\end{corollary}
For $m=5$ the theorem covers every odd $t\ge7$; a single binary boundary
witness at $t=5$ therefore settles that entire fixed-$m$ question too.
It does not settle all $m$.

\begin{proposition}[Odd multiplier closure]\label{prop:multiplier}
If \eqref{eq:target} has a solution at $(m,t)$ and $s$ is any positive
odd integer, it has a solution at $(m,st)$.
\end{proposition}
\begin{proof}
By Lagrange's four-square theorem choose integers $a,b,c,d$ with
$a^2+b^2+c^2+d^2=s$. The integral matrix
\[
 M=\begin{pmatrix}
 a&b&c&d\\-b&a&-d&c\\-c&d&a&-b\\-d&-c&b&a
 \end{pmatrix}
\]
satisfies $M^{\mathsf T}M=sI_4$. Each row sum is odd, since its parity
equals $a+b+c+d\equiv s\equiv1\pmod2$.
Put $B_i=\sum_jM_{ij}A_j$. Then $B_i\equiv J\pmod2$, and expanding
the Hermitian norm sum using real integer coefficients of $M$ gives
\[
 \sum_i B_iB_i^*=s\sum_jA_jA_j^*=4st.
\]
This only uses a constant integral orthogonal design; it does not
assert a four-square composition identity with arbitrary group-ring
Hermitian coefficients.
\end{proof}

\section*{Examples and scope}
At $m=7$, the threshold is $t=13$. Taking $(u,v,w)=(1,1,1)$ gives
$(x_0,x_1,x_2,x_3)=(-1,-2,-2,-2)$, hence constant coefficients
$(-1,-3,-3,-3)$ and all remaining coefficients equal to one.
At $m=9$, the threshold is $21$; choosing $(u,v,w)=(-1,1,1)$ gives
$(-2,-3,-2,-2)$ and constant coefficients $(-3,-5,-3,-3)$.

For $m>3$, $(m^2+3)/4>m$, so the constant-perturbation family cannot
reach the Hadamard boundary. This follows from its sharp lower bound,
not from a failure to search far enough. Applying odd multiplier closure
to any known boundary solution gives infinitely many further energies,
but it cannot lower an energy to the boundary.

The remaining parameter strip for a fixed odd $m$ is finite:
\[
 m\le t<(m^2+3)/4,\qquad t\text{ odd}.
\]
It is an infinite family of strips as $m$ varies. A uniform theorem in
that strip, especially near $t=m$, is still needed. The polynomial
bound above is a proved partial answer to a stronger arithmetic question,
not a new Hadamard-order existence result. Neither the three-square
identity nor this elementary application is claimed as bibliographically
novel without further investigation.

\section*{A lower threshold from skew quadratic characters}
Suppose $m\equiv3\pmod4$ and a polynomial $L\in R_m$ has constant
coefficient zero, all other coefficients signs, and
\[
 L^*=-L,\qquad LL^*=m\cdot1-J.
\]
For every prime $p\equiv3\pmod4$ this is supplied by
$L=\sum_{j=1}^{p-1}\chi(j)z^j$, where $\chi$ is the quadratic character
modulo $p$, as in the Paley character construction~\cite{Paley}.
For completeness, $\chi(-j)=-\chi(j)$ gives skewness. For $s\ne0$,
\[
 \sum_j\chi(j)\chi(j+s)
 =\sum_{j\ne0}\chi(1+s/j)=-1:
\]
as $j$ ranges over the nonzero field elements, $1+s/j$ ranges over all
field elements except one, and a nontrivial character sums to zero.
The constant coefficient is $p-1$, giving the norm identity.

\begin{theorem}[Skew-character large-energy region]\label{thm:skew}
Under those hypotheses on $L$, the norm problem is solvable for every
odd
\[
 t\ge T_m:=\frac{m^2+6m+21}{16}.
\]
In particular this holds for every prime $m\equiv3\pmod4$.
The integer $T_m$ is odd. The threshold is sharp within the family
used in this proof.
\end{theorem}
\begin{proof}
Write $m=4k+3$ and $x=(3-m)/4=-k$. Then
$T_m=k^2+3k+3$ is an odd integer.
For odd integers $a,b,c$, set
\[
 A_1=L+a,\quad A_2=L+b,\quad A_3=L+c,\quad A_4=J+2x.
\]
Every coefficient is odd. Since $L^*=-L$, cross terms with a real
constant vanish, and
\[
 (L+a)(L+a)^*=(m+a^2)\cdot1-J.
\]
The fourth norm is $4x^2\cdot1+(m+4x)J=4x^2\cdot1+3J$.
Thus
\[
 \sum_{i=1}^4 A_iA_i^*
 =\bigl(3m+4x^2+a^2+b^2+c^2\bigr)\cdot1.
\]
The required integer for the last three squares is
$N=4t-3m-4x^2$. It is at least three precisely when $t\ge T_m$.
Also $3m+4x^2=4k^2+12k+9\equiv1\pmod8$, while $4t\equiv4\pmod8$;
hence $N\equiv3\pmod8$. The three-square theorem supplies odd
$a,b,c$ with squared sum $N$. Conversely, any three odd integers have
squared sum at least three, proving sharpness in this family.
\end{proof}

At $m=7$ this threshold is seven, so the entire broader problem is
settled for that fixed length as well. At $m=11$ it is thirteen, leaving
only the binary boundary case $t=11$ below the threshold. For large
primes the threshold remains quadratic in $m$, so it does not reach
the Hadamard boundary uniformly. The prime-power version of Paley's
character construction lives in the group ring of the additive field,
which need not be cyclic; no cyclic prime-power extension is asserted.

\begin{proposition}[A zero-aligned character ansatz is impossible]
With the same $m,L$, it is impossible that all four elements have the
form
\[
 A_i=u_iL+v_i+2w_iJ,\qquad u_i,v_i\text{ odd},\quad w_i\in\Z,
\]
and their Hermitian norm sum is a scalar multiple of $1$.
\end{proposition}
\begin{proof}
The coefficient multiplying $J$ in the sum of norms is
\[
 K=-\sum_i u_i^2+4\sum_i(w_iv_i+mw_i^2).
\]
Every summand in the last parentheses is even, since its parity is
$w_i+w_i^2=0$. The four odd squares sum to four modulo eight.
Therefore $K\equiv4\pmod8$, so it cannot vanish. The elements $J$ and
$1$ are linearly independent because $m>1$.
\end{proof}
This is an obstruction to a restricted character ansatz, not to the
general norm problem or to Hadamard existence. Theorem~\ref{thm:skew}
avoids it by using an odd coefficient of $J$ in the fourth element.

\section*{Exact checks}
The standard-library script \path{verify_large_energy.py} constructs and
checks examples at the threshold and several larger odd energies for
every odd $m\le101$. It also tests odd multiplier closure using a direct
cyclic autocorrelation check. The general assertions rest on the paper
proofs; finite tests are additional checks only.

% Self-contained mathematical fragment. Uses proposition/proof environments,
% standard amsmath/amssymb, hyperref, and the parent note's \Z and \R macros.
\section{Moment-semigroup progress: which moves can preserve positivity?}
All results in this section have elementary proofs. They do not establish
new Hadamard orders or a general positivity-conversion theorem. Their
purpose is to discriminate between possible algebraic-statistics tools.
In particular, generating the integer kernel is much weaker than connecting
its nonnegative fibers.

\subsection{Cube moves are frozen at every saturated Hadamard design}
An observation is a bit vector in $\{0,1\}^d$, with integer multiplicity.
A coordinate-subcube parity move fixes coordinates outside a set $T$ and
adds one to the vertices of one parity inside that subcube while subtracting
one from the other parity. When $|T|\ge3$, the move preserves every
quadratic moment. Both signs of every translated move are allowed.

\begin{proposition}[Isolation under all coordinate-cube moves]
Let a binary strength-two orthogonal array have $N$ rows and $d=N-1$
columns, with $N\ge8$. No nonzero integer multiple of any individual
coordinate-subcube parity move of dimension at least three is feasible
from its multiplicity vector while preserving nonnegativity.
\end{proposition}
\begin{proof}
Replace bits by signs and adjoin the all-one column. The resulting square
matrix $H$ has $H^{\mathsf T}H=NI$, hence $HH^{\mathsf T}=NI$. Its
observations are distinct, and every two have Hamming distance $N/2$:
their sign inner product after deleting the common first coordinate is
$-1$, so $d-2\operatorname{dist}(x,y)=-1$. In particular no two
supported observations are at Hamming distance two.

Each parity class of a coordinate subcube of dimension at least three
contains a pair of vertices at Hamming distance two: flip any two of its
free coordinates. Either orientation of the move subtracts from all
vertices of one such class. That class cannot lie in the support of the
array. Some subtracted coordinate therefore has multiplicity zero, and
every positive integer multiple makes it negative.
\end{proof}

Flipping a single coordinate produces a different array in the same
quadratic-moment fiber: unbiasedness and pairwise independence are
preserved, and none of the new distance-one neighbors was supported
before. Thus even the collection of all translated coordinate-cube moves,
in all dimensions, cannot connect these nonnegative fibers and cannot
be a Markov basis for them. This is not just a failure of a low-dimensional truncation of the
origin-anchored kernel basis. General non-coordinate trades are necessary.

\paragraph{An exact degree-four trade outside this move family.}
Take the Walsh matrix $H_{ab}=(-1)^{\langle a,b\rangle}$ of order eight,
with $a,b\in\{0,1\}^3$ ordered by their integer values. Replace its
first four rows by
\[
 r_i'=\frac12(r_0+r_1+r_2+r_3)-r_i\qquad(0\le i<4).
\]
The row transformation $J_4/2-I_4$ is orthogonal. Each column in these
four rows has product $+1$, so its sum is $0$ or $\pm4$; the transformed
entries are again signs. The first column remains all ones. Consequently
both old and new matrices are Hadamard, and their seven-column arrays
have the same quadratic moments.

Represent a bit pattern by the mask whose bit $j$ is its $j$th coordinate.
The old array has masks
\[
 \{0,85,102,51,120,45,30,75\}.
\]
The trade removes $\{0,85,102,51\}$ and inserts
$\{119,34,17,68\}$. It preserves nonnegativity and has degree four
(four observations on either side). Both endpoints are isolated under
every coordinate-cube move by the proposition. The supplied checker
enumerates all $939$ coordinate subcubes of dimensions three through
seven and checks both orientations at both endpoints. This is the
classical kind of closed-quadruple switching, independently checked here,
not a claim of a new switching construction.

\subsection{A second obstruction: move order in the anchored basis}
\begin{proposition}[Sharp maximum-weight constraint]
Every unbiased pairwise-independent distribution on $\{0,1\}^d$ whose
support has Hamming weight at most $K$ satisfies
\[
 K\ge\left\lceil\frac{d+1}{2}\right\rceil.
\]
The bound is attained by a positive integer orthogonal array for every
$d\ge2$.
\end{proposition}
\begin{proof}
For the Hamming weight $W=\sum_i Z_i$, pairwise independence gives
$\mathbb EW=d/2$ and $\mathbb EW^2=d(d+1)/4$. Since $W^2\le KW$,
the asserted lower bound follows.

For attainment put $s=\lfloor d/2\rfloor$ and
$C_s=\binom{2s}{s}/(s+1)$. If $d=2s+1$, take every weight-$(s+1)$
vector once and $C_s$ copies of the zero vector. If $d=2s$, take every
weight-$s$ and every weight-$(s+1)$ vector once, together with $C_s$
zero vectors. In either case the run count is
\[
 N=2\binom{2s}{s},
\]
and a fixed coordinate appears in $N/2$ rows while any fixed pair
appears in $N/4$ rows. For odd $d$, these identities are
\[
 \binom{2s+1}{s+1}+C_s=2\binom{2s}{s},\quad
 \binom{2s}{s}=N/2,\quad
 \binom{2s-1}{s-1}=N/4.
\]
The even case is obtained by deleting a coordinate from the odd one,
or directly by Pascal's identity. The maximum weight is exactly
$s+1=\lceil(d+1)/2\rceil$.
\end{proof}

Start from the canonical signed solution supported at weights zero,
one and two. If only origin-anchored basis moves $M_T$ with $|T|\le K$
are used, with arbitrary integer coefficients, the result is still
supported at weights at most $\max\{2,K\}$. Hence conversion to a positive
design requires a move order at least $\lceil(d+1)/2\rceil$ when this
exceeds two. This claim concerns that anchored basis. It does not rule
out translated three-dimensional moves based at large-weight patterns;
the preceding isolation proposition is a separate statement about their
behavior at the saturated endpoints.

\begin{proposition}[Exact limitation of monotone triple conversion]
Write $L=N/4>0$. Starting from the canonical signed solution, allow only
nonnegative multiples of the origin-anchored three-cube moves which add
to singletons and triples and subtract from the empty set and pairs.
For $d=4$ conversion is possible exactly when $L\ge2$; for $d=5$ it is
possible exactly when $3\mid L$; for $d\ge6$ it is impossible.
\end{proposition}
\begin{proof}
Let $t_T\ge0$ be the multiplicity of each triple move, and let
$T_0=\sum_{|T|=3}t_T$. The final singleton counts require
$3T_0\ge dL(d-3)$, while pair nonnegativity requires
$3T_0\le\binom d2L$. For $d\ge6$ these inequalities are inconsistent.

For $d=5$ they force equality and $T_0=10L/3$, hence $3\mid L$.
Conversely, setting every triple multiplicity to $L/3$ works: the only
remaining positive counts are $2L/3$ at the empty pattern and $L/3$
at each of the ten triples. For $d=4$, the case $L=1$ would give five
orthogonal nonzero columns in $\R^4$, which is impossible. For every
$L\ge2$, set all four triple multiplicities to $\lfloor L/2\rfloor$.
The resulting empty, singleton, pair and triple counts are respectively
$2L-4\lfloor L/2\rfloor$, $-L+3\lfloor L/2\rfloor$,
$L-2\lfloor L/2\rfloor$, and $\lfloor L/2\rfloor$, all nonnegative.
\end{proof}
For instance the unrestricted eight-run design on five coordinates
exists by truncating Walsh order eight, but the triple-only conversion
above cannot reach any positive design at $(N,d)=(8,5)$. The obstruction
belongs to the move strategy, not to the moment target.

\subsection{Why averaging over all coordinate permutations loses the target}
\begin{proposition}[Exponential cost of full permutation symmetry]
Let $N_{\mathrm{sym}}(d)$ be the least positive run count for a binary
strength-two array whose pattern multiplicities are invariant under all
coordinate permutations. Put
\[
 k_d=\min\{k\in\{0,\ldots,d\}:(2k-d)^2\le d\}.
\]
For $d\ge2$,
\[
 \binom d{k_d}\ \le N_{\mathrm{sym}}(d)
 \le 2\binom{2\lfloor d/2\rfloor}{\lfloor d/2\rfloor}.
\]
In particular $N_{\mathrm{sym}}(d)=\Theta(2^d/\sqrt d)$.
\end{proposition}
\begin{proof}
The Hamming weight has mean $d/2$ and variance $d/4$. At least one
occupied weight $k$ has $(2k-d)^2\le d$. Full permutation symmetry
then forces each of the $\binom dk$ patterns at that weight to have
positive integer multiplicity. Throughout this central interval,
$\binom dk\ge\binom d{k_d}$, proving the lower bound. The Catalan-shell
construction above proves the upper bound. Finally
$k_d=d/2-\sqrt d/2+O(1)$, so Stirling's formula gives the matching
orders. More precisely the lower binomial is asymptotic to
$e^{-1/2}2^d\sqrt{2/(\pi d)}$.
\end{proof}
Thus, for sufficiently large $d$, exact global symmetrization cannot
preserve the saturated run count $N=d+1$. Selective symmetry or symmetry
breaking is essential in that regime. The
Catalan-shell constructions are related to the recent paper
\emph{Catalan Bounds for Symmetric Strength-Two Orthogonal Arrays}
by R. C. Karunanayaka,
\url{https://doi.org/10.1002/jcd.70039}; no priority is claimed for them.
The bound above concerns growth order and maximum weight, not the exact
minimum run-count conjecture in that paper.

\subsection{A positive result: a quadratic bound for the fixed-dimension conductor}
\begin{proposition}[An explicit cofinite central ray]
For each fixed $d\ge2$, every sufficiently large $N\equiv0\pmod4$
admits a binary strength-two array on $d$ coordinates. An explicit
threshold is
\[
 N\ge C_d=4(A-1)(B-1)<3(d+1)^2,
\]
where $a=4A$ is the least power of two at least $d+1$, and $b=4B$ is
the first order at least $d+1$ in the sequence
\[
 b_e=\begin{cases}3^e+1&e\text{ odd},\\2(3^e+1)&e\text{ even}.
 \end{cases}
\]
Here $N$ is positive, also when $C_d=0$.
\end{proposition}
\begin{proof}
Sylvester gives order $a$. Paley's two constructions give order $b$:
odd $e$ uses type I and even $e$ uses type II. Every $b_e$ has exact
two-adic valuation two, while $A$ is a power of two, so $\gcd(A,B)=1$.
Truncating the normalized matrices gives arrays with $d$ coordinates
and run counts $a,b$. Concatenating observations adds moment vectors.
The elementary two-generator numerical-semigroup theorem therefore
gives every run count $4M$ with $M\ge(A-1)(B-1)$. For completeness,
choose $u\in\{0,\ldots,B-1\}$ with $uA\equiv M\pmod B$. Then
$M-uA\ge-B+1$, so its integer quotient by $B$ is nonnegative.

We have $a<2(d+1)$. Consecutive displayed $b_e$ grow by a factor less
than six, giving $b<6(d+1)$, including the initial case $b=4$.
Consequently $C_d<ab/4<3(d+1)^2$.
\end{proof}
This result uses only classical Sylvester and Paley constructions, not
the current project's asymptotic theorem. It proves finiteness and an
explicit bound for central-ray holes when $d$ is fixed. It remains far
from the diagonal $N=d+1$ needed for Hadamard existence. It is an
elementary application of known constructions, not a claim of a new
Hadamard bound.

\paragraph{Reproduction.}
The standard-library script \path{moments/check.py} checks the isolated
endpoints and non-cube trade, all $939$ coordinate subcubes, the shell
moment identities through $d=40$, low-dimensional explicit arrays,
triple-conversion witnesses, and sample decompositions above the
conductor thresholds. Results are in \path{moments/certificates.json}.
None of these checks is a Lean formalization or a proof of unrestricted
Hadamard existence.
% Standalone section fragment. Uses amsmath, amssymb, amsthm and hyperref.
% These elementary reductions are not claimed to have new bibliographic priority.
\section{A sharper obstruction to bounded row-exchange augmentation}

This pass gives a concrete negative result: replacing two rows is not
enough for every augmenting exchange. The following finite certificate
raises the necessary universal deletion bound from two to at least three.
It does not establish any universal upper bound.

\subsection{A small-dimensional exact enumeration theorem}
\begin{proposition}[Cube intersections through information coordinates]
Let $V\subseteq\mathbb R^n$ be a $q$-dimensional linear subspace. Then
$|V\cap\{\pm1\}^n|\le2^q$. There is a set of $q$ coordinate positions
such that the projection to those positions is injective on $V$.
Consequently, if $V$ is given by rational linear equations, all its
sign vectors can be enumerated by testing at most $2^q$ assignments,
using exact rational arithmetic. Antipodal normalization reduces this
to $2^{q-1}$ assignments when $q>0$.
\end{proposition}
\begin{proof}
Choose a basis matrix for $V$ and $q$ linearly independent coordinate
rows. Projection to those coordinates is an isomorphism onto
$\mathbb R^q$. A sign vector projects to a sign vector, so there are
at most $2^q$ possibilities. Each choice determines a unique candidate
in $V$; check whether all remaining coordinates are signs. For rational
data, elimination and all checks are exact. Negation pairs all nonzero
cube vectors, allowing one information coordinate to be fixed positive.
\end{proof}

For a partial matrix with $m$ orthogonal rows at order $n$, deleting $k$
rows leaves a kernel of dimension $q=n-m+k$. The enumeration cost is
therefore exponential in the defect plus the deletion count, not in
the full order. After enumeration, a $(k+1)$-clique in the candidate
orthogonality graph is exactly an augmenting exchange for that deletion
set. This gives an exact decision procedure for a specified exchange,
not a theorem that such an exchange exists.

There is also a rigid equality case, useful when classifying residual
frames rather than only counting candidates.
\begin{proposition}[Maximum cube intersections are duplication spaces]
The bound $|V\cap\{\pm1\}^n|=2^q$ holds if and only if, after
permuting and negating coordinates, $V$ consists of vectors constant
on each of $q$ nonempty coordinate blocks.
If $V$ is the residual space of a partial Hadamard matrix, all these
blocks have size $n/q$. In that case completion is equivalent to the
existence of a Hadamard matrix of order $q$.
\end{proposition}
\begin{proof}
Choose information coordinates as above. Equality means every sign
assignment to them extends to a sign vector. Each other coordinate is
a linear form $f(\varepsilon)=\sum_i a_i\varepsilon_i$ taking only
values $\pm1$ on the full $q$-cube. Its maximum is
$\sum_i|a_i|=1$, whereas averaging $f^2$ gives $\sum_i a_i^2=1$.
These two equalities force at most one nonzero coefficient, which is
$+1$ or $-1$. Thus every coordinate is a signed copy of an information
coordinate. The converse is immediate.

For a partial Hadamard residual space, the orthogonal projector has
constant diagonal $q/n$. The projector on a duplication block of size
$s$ has diagonal $1/s$, so $s=n/q$ for every block. Inner products of
sign vectors are then $(n/q)$ times the inner products of their $q$
block values. A full orthogonal basis of sign vectors exists precisely
when those block values form an order-$q$ Hadamard matrix.
\end{proof}

This shows why even a maximal number of sign candidates is not, without
additional structure, an automatic completion theorem: the smaller
Hadamard existence problem appears explicitly in the equality case.

\subsection{An exact order-twenty-four certificate}
Encode a sign row $x$ by the integer whose bit $j$ is one exactly when
$x_j=-1$; the least significant bit is coordinate zero. The sixteen
rows in this example are
\begin{align*}
 &0,\ 16773120,\ 16519104,\ 10359454,\ 7684730,\ 5820908,\ 9389414,\ 1547504,\\
 &6736278,\ 11838780,\ 3701846,\ 13790378,\ 10646512,\ 2963882,\ 12892774,\ 5974836.
\end{align*}
Each row has length $24$. All distinct row pairs differ in exactly
twelve positions, so their dot products vanish.

\begin{proposition}[A three-row deletion barrier]
This partial matrix has no augmenting exchange that deletes at most
two rows. Deleting three rows does admit an augmenting exchange and
a subsequent full completion. Thus any universal row-exchange
augmentation theorem must permit at least three deletions.
\end{proposition}
\begin{proof}[Exact finite certificate]
Independently of the low-dimensional enumeration procedure, the
verifier exhausts all $2^{23}=8388608$ sign vectors normalized to have
coordinate zero positive. For each vector, record the set of original
rows to which it is \emph{not} orthogonal. Only vectors with at most
two such rows can participate after at most two deletions. There are
exactly $56$: sixteen have a singleton nonorthogonality support, and
forty have support of size two. No vector has empty support.

For each of the sixteen single-row deletions there is exactly one
normalized candidate. Of the $\binom{16}{2}=120$ two-row deletions,
one hundred have exactly two candidates, and twenty have four
candidates. None of the latter candidate sets contains three mutually
orthogonal rows. This exhausts all $1+16+120=137$ deletion sets and
rules out an augmentation deleting at most two rows.

Delete rows with zero-based indices $0,5,6$. Insert the four rows
encoded by
\[
 353924,\quad1104648,\quad2786850,\quad4759768.
\]
The thirteen retained rows and these four new rows are pairwise
orthogonal. The resulting seventeen rows complete by adding
\[
 5133676,\ 7583694,\ 9432922,\ 10076646,\ 10684492,\ 12653842,\ 15326908.
\]
The verifier checks the full order-$24$ Gram matrix by exact Hamming
distances. The positive certificate and the exhaustive negative
certificate together prove that the minimal deletion count is three.
\end{proof}

This is a counterexample to universal two-row augmentation, not a
counterexample to Hadamard existence. It also demonstrates that
``single-row frozen'' understates the actual obstruction: some
two-row deletions have genuinely new candidates, but their mutual
orthogonality graph still prevents an augmentation.

The matrix was obtained from the first terminal case in a fixed-seed
thirty-run order-$24$ experiment with a canonical three-row prefix.
The row-count histogram was nineteen terminal sixteen-row matrices,
two terminal twelve-row matrices, and nine full matrices. Those
frequencies describe the sampling algorithm only. No claim about all
partial matrices or all equivalence classes is inferred from them.

\subsection{Files and independent reproduction}
All files for this pass are in \path{research/full_conjecture/exchanges/}.
Run \path{sample.py}, \path{augment.py}, then \path{verify.py}.
The first two scripts use exact free-coordinate enumeration and clique
search. The independent verifier does not use that enumeration or
rational elimination: it traverses the entire normalized sign cube,
records nonorthogonality supports, and checks every small deletion set.
Its report is \path{verification.json}; the complete compact list of
relevant candidate supports is in \path{independent24.json}.
No Lean claim or bibliographic novelty claim is made.
\section{A quartic obstruction to a benign optimization landscape}
The previous note ruled out a naive local-maximization argument for
the entrywise one-norm. We now test the smoother alternative
\[
 Q(U)=\sum_{ij}u_{ij}^4\qquad(U\in O(n)).
\]
It satisfies $Q(U)\ge1$, with equality exactly when $\sqrt nU$ is a
Hadamard matrix. The arbitrary-exponent framework is classical
\cite{ArbitraryP}; the calculations below are independently derived,
and no novelty claim is made.

\begin{proposition}[An infinite family of strict quartic traps]
Let $A=2J_5/5-I_5$, and let $B=H/\sqrt k$ for any Hadamard matrix
$H$ of order $k$. Then $U=A\otimes B$ is a strict local minimum of
$Q$ on $O(5k)$, with value $29/25$. In fact its Hessian obeys
\[
 \left.\frac{d^2}{dt^2}Q(Ue^{tX})\right|_{t=0}
 \ge\frac4{5k}\|X\|_F^2
 \qquad(X^{\mathsf T}=-X).
\]
For $k=4\cdot2^a$, these are nonglobal local minima at the admissible
orders $20\cdot2^a$.
\end{proposition}
\begin{proof}
For an orthogonal matrix $V$, entrywise differentiation gives
\begin{align*}
 \left.\frac{d}{dt}Q(Ve^{tX})\right|_{0}
 &=4\operatorname{tr}((V^{\circ3})^{\mathsf T}VX),\\
 \left.\frac{d^2}{dt^2}Q(Ve^{tX})\right|_{0}
 &=12\sum_{ij}v_{ij}^2(VX)_{ij}^2
   +4\operatorname{tr}((V^{\circ3})^{\mathsf T}VX^2).
\end{align*}
At the claimed $U$, the symmetric matrix
\[
 M=U^{\mathsf T}U^{\circ3}
  =\left(\frac7{25}I_5-\frac6{125}J_5\right)\otimes\frac1kI_k
\]
has eigenvalues $1/(25k)$ and $7/(25k)$. Thus $U$ is stationary.
Its smallest squared entry is $4/(25k)$. Since
$\|UX\|_F=\|X\|_F$ and $X^2=-X^{\mathsf T}X$,
\[
 Q''(0)\ge\left(\frac{48}{25k}-\frac{28}{25k}\right)\|X\|_F^2
 =\frac4{5k}\|X\|_F^2.
\]
The positive definite tangent Hessian proves strict local minimality.
Direct summation gives $Q(A)=29/25$ and $Q(B)=1$, hence the stated
value. A Hadamard matrix of order $20$ is furnished by the certified
length-five Goethals--Seidel witness; tensoring it with the Sylvester
matrix of order $2^a$ gives a point of value one at the same order.
Consequently these local minima are not global.
\end{proof}

The obstruction is uniform, not an inference from numerical Hessian
eigenvalues. The exact script checks the algebra at orders $20,40,80$.
It shows that changing the one-norm objective to the fourth-power
objective does not by itself remove the false-minimum problem.

\begin{proposition}[Coordinate rotations can miss a descent direction]
At order $28$, let
$U=(2J_7/7-I_7)\otimes(H_4/2)$. This matrix is stationary for $Q$.
Every single Givens rotation has strictly positive second derivative,
but a joint tangent direction has negative second derivative.
\end{proposition}
\begin{proof}
For $A_n=2J_n/n-I_n$, direct substitution in the Hessian formula yields
\[
 \left.\frac{d^2}{dt^2}Q(A_ne^{tX})\right|_0
 =-\frac{4(n-6)}n\|X\|_F^2
  +\frac{8(n^2-16)}{n^3}\|X\one\|^2.
\]
Take $n=7$. A single Givens direction in the tensor matrix has
curvature $34/343$ when its two coordinates belong to distinct
order-seven blocks, and $412/343$ when they belong to the same block.
Both are positive. Set
$X_0=E_{01}+E_{12}-E_{02}$, where
$E_{ij}=e_ie_j^{\mathsf T}-e_je_i^{\mathsf T}$. Then
$X_0\one=0$ and $\|X_0\|_F^2=6$. The joint direction
$X=X_0\otimes I_4$ gives curvature $-24/7$, since the fourth-power
objective of the normalized Hadamard tensor factor is one.
All values are also checked by the exact full tangent Hessian script.
\end{proof}
Thus testing all coordinate-pair second derivatives is not a sufficient
second-order test even at an admissible order. This matters when
importing coordinate descent or block-rotation methods from numerical
optimization. The first proposition supplies genuine strict minima;
the second supplies a different failure, namely a saddle hidden from
every individual coordinate rotation.

\section{Consequences for the next research step}
The arithmetic result replaces an unrestricted search over odd
coefficients by explicit uniform constructions in a large region.
Its remaining task is to lower the threshold to the minimum energy,
not to multiply known solutions by larger odd scalars. In the
skew-character construction one channel still has a large row sum;
the all-four aligned-character shortcut is ruled out by a parity
obstruction. A successful extension must change that structure.

The moment results require trades beyond coordinate subcubes and a
positivity argument adapted to sparse saturated fibers. Their
fixed-dimension conductor bound does not control the growing diagonal
$N=d+1$. The exchange certificate raises the necessary universal
deletion allowance to three but establishes no universal upper bound.
The quartic result means that a proposed optimization proof needs a
global argument or a way to escape genuine strict local minima.

No claimed result here proves new universal Hadamard existence. The
general proof remains open in this project. The deductions have been
reviewed for their stated domains, and exact computations are reproduced
from standard-library Python scripts. The scripts are checks of finite
certificates; the infinite assertions rest on the proofs in this note.

\section{Files and verification}
Branch sources are under \path{arithmetic/}, \path{moments/},
\path{exchanges/}, and \path{geometry/}, relative to
\path{research/full_conjecture/}. Run
\path{arithmetic/verify_large_energy.py}, \path{moments/check.py},
\path{exchanges/test_cube.py}, \path{exchanges/verify.py}, and
\path{geometry/verify.py}. The full-cube exchange verifier checks all
$2^{23}$ normalized sign vectors independently of the rational
nullspace search. The geometry proof uses an analytic Hessian lower
bound, not floating-point eigenvalue tests. No external research
proof code is imported, and no Lean work or Git commits were made.

\begin{thebibliography}{9}
\small\raggedright

\bibitem{Ankeny} N. C. Ankeny, \emph{Sums of three squares},
Proc. Amer. Math. Soc. 8 (1957), 316--319, Theorem 1.
\url{https://doi.org/10.1090/S0002-9939-1957-0085275-8}.
\bibitem{Paley} R. E. A. C. Paley, \emph{On orthogonal matrices},
J. Math. Phys. 12 (1933), 311--320.
\url{https://doi.org/10.1002/sapm1933121311}.

\bibitem{ArbitraryP} T. Banica and I. Nechita,
\emph{Almost Hadamard matrices: the case of arbitrary exponents},
Discrete Appl. Math. 161 (2013), 2367--2379.
\url{https://arxiv.org/abs/1211.2669}.
\end{thebibliography}
\end{document}
